% Literature review (B. Horvath's condensed draft, pasted verbatim; [n] -> \cite keys).
% Cut: the final paragraph of "Why DLV and SANOS, and where GROMM fits" (numerical measure change for 48 options),
% which contradicts GROMM's explicit risk-neutral measure. See CUTS.md.
\section{Literature Review}
\label{app:literature-review}

Modelling joint stock--option dynamics requires fitting market quotes, preserving static arbitrage
freedom across strikes and maturities, and excluding dynamic trading arbitrage. These are
distinct objectives: a sequence of admissible surfaces need not admit a common equivalent
martingale measure. We organise the literature by its surface representation, or codebook, and
the arbitrage restrictions imposed on its evolution.

\paragraph{Implied-volatility dynamics and dynamic consistency.} Early empirical studies identified
low-dimensional patterns in implied-volatility movements~\cite{SkiadopoulosHodgesClewlow2000,ContDaFonseca2002Dynamics,FenglerHardleVilla2003}. Stochastic surface models and
associated drift restrictions were investigated in~\cite{ContVol}, while subsequent factor approaches developed
their statistical modelling~\cite{Hafner2004,FenglerHardleMammen2007}. Ordinary implied-volatility factors, however, do not automatically
preserve admissible option prices, and fitting historical dynamics does not establish a joint pricing
measure. This distinction remains relevant to modern latent-variable models: dimensionality
reduction alone provides neither a statically admissible decoder nor dynamically consistent
evolution.

Direct market models address dynamic consistency explicitly. No-arbitrage restrictions for
stochastic implied volatilities, including a fixed-strike term structure, are developed in~\cite{Schonbucher1999}; drift
and existence results for the fixed-strike problem follow in~\cite{SchweizerWisselTerm2008}. The complementary multi-strike
construction~\cite{SchweizerWissel2008} replaces ordinary implied volatilities at one maturity by a price level and a field
of local implied volatilities. These encode static restrictions across strikes; separate coefficient
restrictions ensure dynamic consistency. This motivates selecting admissible coordinates before
prescribing their evolution. Local implied volatility here differs from Dupire local volatility, so
our method is conceptually related rather than a discretisation of those equations.

The symmetric SSVI specification in~\cite{DysLitElAmrani2021} cannot reconcile its martingale requirements with
vanishing total implied variance at expiry in the desired nontrivial model. It instead leads to
implied-volatility bubbles with a restricted arbitrage-free lifespan. This model-specific impossibility
illustrates why static smile admissibility does not ensure dynamic consistency through
expiry.

\paragraph{Local-volatility and alternative codebooks.} The local-volatility market codebook in~\cite{CarmonaNadtochiy2009}
is the closest continuous ancestor of our approach: local volatility represents the current option
surface, while consistency conditions govern its joint evolution with the underlying. Alternative
codebooks use tangent L\'evy characteristics~\cite{TangentLevy} or time-inhomogeneous L\'evy exponents in an
HJM framework~\cite{KallsenKruhner2015}. The latter supplies necessary and sufficient no-arbitrage conditions and
constructive existence and uniqueness results. Empirical estimation and simulation of tangent
L\'evy models are developed in~\cite{CarmonaMaNadtochiy2017}. Thus dynamically consistent surface models predate neural
simulators; the relevant distinctions concern representation, computation and statistical flexibility.

Consistent recalibration~\cite{DysLitGambaraTeichmann2021} uses forward characteristics and generalised Hull--White extensions,
illustrated through an extended Bates model and neural calibration. Parameters evolve while a
compensating extension preserves the current surface, subject to admissibility and consistency
conditions. This shares the preceding market-model lineage but evolves an extended pricing-model
representation rather than DLV coordinates. Its compensation mechanism is substantive:
recalibrating a conventional pricing model independently at each date would not itself produce a
consistent stochastic market model. Analytical consistency must be distinguished from numerical
recalibration accuracy: the theoretical framework does not certify every approximate neural
inverse or arbitrary proposed parameter reset.

\paragraph{Learned models addressing dynamic arbitrage.} Affine factors of normalised call prices
and neural SDEs are combined in~\cite{Cohen02012023}. The price-factor codebook differs from DLV, although
both seek to preserve admissible prices under learned dynamics. An invariant factor polytope
structurally preserves finite-lattice static inequalities. An HJM-type dynamic condition is also
derived, but the rolling-lattice implementation pursues it through factor selection and residual
diagnostics rather than exact enforcement for all fixed contracts. The theoretical condition,
implemented static guarantee and numerical dynamic criterion therefore have different scopes.

A complementary approach learns a change of measure on historical simulator paths, targeting
zero conditional trading gains~\cite{removedrift}. DYSANOS~\cite{DYSANOS2026} applies this to a smooth SANOS generator,
combining historical dynamics, structural static admissibility and a numerical risk-neutral density
for selected instruments. It consequently belongs among models explicitly addressing dynamic
arbitrage, while its finite-universe validation differs from a theorem providing one pricing measure
for the full option continuum. Such a measure change also inherits the surface representation
and its pathwise static properties from the input simulator.

\paragraph{Structural static admissibility.} Other generators preserve static admissibility without
establishing joint risk-neutral dynamics. DLV-based simulators~\cite{wiese2019simulate,wiese2021multiassetspotoptionmarket} enforce admissibility on
their native lattice. Path-dependent SSVI dynamics~\cite{AndresBoumezouedJourdain2023} evolve within an admissible parametric
family. These constructions preserve static restrictions at simulated dates, but that property alone
does not establish dynamic no-arbitrage. Their codebooks also differ: the former parameterise
discrete local variances, whereas the latter restrict entire surfaces through SSVI parameters.

The VAE literature needs a separate distinction. The latent surface representation in~\cite{BergeronEtAl2021} has no
hard static guarantee. By contrast,~\cite{vae_ning2023arbitrage} decodes into parameters of arbitrage-free regime-switching
or L\'evy additive pricing models. Admissibility comes from the pricing representation, not an
arbitrage penalty on ordinary implied-volatility outputs. Nevertheless, sampling pricing-model
parameters across dates does not establish dynamically consistent traded-price evolution. A
statically admissible snapshot generator should therefore not be assessed as though it had already
constructed a multi-step market model.

\paragraph{Soft constraints and unconstrained generators.} Static restrictions are instead encouraged
through losses, reweighting or output corrections in~\cite{ZhangLiZhang2021,ContVuletic2023,ChoudharyJaimungalBergeron2024,VuleticCont2024}. FuNVol combines functional principal
components and neural SDEs for joint asset--surface movements, with penalties controlling static
violations. VolGAN and related reweighting methods also seek empirically admissible outputs
without restricting every decoded state to an admissible family. Low violation rates do not
guarantee admissibility between diagnostic points or beyond the training distribution. The original
VAE~\cite{BergeronEtAl2021} is more accurately classified as unconstrained than as a guaranteed or necessarily
penalty-enforced construction.

Recent diffusion and latent-flow approaches~\cite{JinAgarwal2025,BrooksEtAl2026,HaoJi2026,HanEtAl2026,YangWangZhangLi2026} extend these architectures, but their arbitrage
control remains approximate in the implementations surveyed. Their statistical performance,
static constraint violations and dynamic consistency require separate evaluation; snapshot generation,
one-step forecasting and multi-step simulation are likewise different tasks. This classification
concerns the reported constructions and does not suggest that a generative architecture itself
prevents structural arbitrage constraints from being incorporated.

\paragraph{Related smile-dynamics literature.} Forward-variance models~\cite{DysLitBergomi2005,bergomi2008smile} and stochastic-volatility
smile expansions~\cite{bergomi2011smile} investigate the volatility structures underlying smile dynamics. Leverage-based
analysis~\cite{DysLitCilibertiBouchaudPotters2008} explains movements intermediate between sticky-strike and sticky-delta rules.
These provide economic targets for surface simulators, rather than the same nonparametric
surface codebook. They are particularly relevant when evaluating whether a simulator reproduces
spot--surface dependence beyond matching individual surface distributions.

Rough-SABR and asymptotically arbitrage-free approximation results~\cite{DysLitFukasawaGatheral2022,DysLitFukasawa2022} connect implied-volatility
approximations to pricing-model drift restrictions. Their asymptotic guarantees differ
from exact admissibility of generated full surfaces. Normalising volatility transforms and smile
parametrisation~\cite{lucic2021normalizing} provide further background on admissible coordinates, rather than an
established joint dynamically arbitrage-free surface simulator. These adjacent contributions
clarify which statistical and geometric properties a codebook should preserve without conflating
approximation theory with a market-model existence theorem.

\paragraph{Why DLV and SANOS.} Our choice follows the local-volatility
branch. Implicit forward-equation pricing~\cite{AndreasenHuge} and DLV~\cite{DLV} replace coupled inequalities on
implied volatilities with constrained local-variance coordinates and a structured pricing map.
Local volatility serves as a surface coordinate system: it does not force historical stock--surface
dynamics to follow a fixed local-volatility diffusion. This separates learned surface evolution from
structural static admissibility.

The computational attraction is especially relevant for repeated decoding during optimisation:
solving a fine forward PDE for every candidate surface would be expensive. Positivity constraints
on codebook coordinates are also easier to impose than the coupled strike and maturity restrictions
on directly generated implied volatilities.

SANOS~\cite{SANOS2026} provides the smooth admissible extension needed to evaluate contracts with fixed
cash strikes and calendar expiries beyond the native lattice. Transformed coordinates permit unconstrained latent learning:
\[
h_t\in\mathbb{R}^d\ \longmapsto\ u_t=\Phi(h_t)\ \longmapsto\ C_t(\tau,k)=D_{\mathrm{SANOS}}(u_t;\tau,k).
\]
The codebook is the DLV-derived SANOS representation, not the latent vector or transition alone.
Its lineage runs from local-volatility market codebooks through computational DLV to
smooth SANOS decoding. The connection to~\cite{SchweizerWissel2008} is the principle of choosing an admissible
parametrisation before specifying dynamics, rather than identical coordinate formulas.
